\documentclass[12pt]{article}
\usepackage[utf8]{inputenc}
\usepackage[T1]{fontenc}
\usepackage{hyperref}
\title{San Baltazar Loxicha Zapotec Language Resources for K-12 Curriculum Development}
\author{Compiled by Claude Sonnet 5.5 (claude-sonnet-5-5) for Andr\'as Kornai}
\date{}
\begin{document}
\maketitle

\section{Introduction}
San Baltazar Loxicha Zapotec (ISO 639-3 zpx, Glottocode sanb1241) is an Otomanguean language of southern Oaxaca, Mexico. It is spoken in San Baltazar Loxicha and Santa Catarina Loxicha. The speaker figure usually cited is about 1,500, and it comes from an old census. Very little published material exists for this variety. No school materials, primers or dictionaries turned up in the searches.

\section{Primary Educational Resources}
No primers, textbooks or school programmes for this language were found in search results.

\section{Linguistic and Reference Materials}
\begin{itemize}
\item Wikipedia overview of the language: \url{https://en.wikipedia.org/wiki/San_Balt%C3%A1zar_Loxicha_Zapotec}
\item Glottolog entry (classification and bibliography pointers): \url{https://glottolog.org/resource/languoid/id/sanb1241}
\item Terrence Kaufman, \emph{Field notes on San Balt\'azar Loxicha Zapoteco} (1997; elicitation and field notes; PDF and TIFF scans), held in the AILLA archive at the University of Texas: \url{https://islandora-ailla.lib.utexas.edu/islandora/object/ailla%3A135184}
\end{itemize}

\section{Religious and Community Resources}
No community organisations or scripture translations were found. Global Recordings Network lists Christian audio material for this language (see below).

\section{Digital and Audio Resources}
\begin{itemize}
\item Global Recordings Network, Zapotec, San Baltazar Loxicha: evangelism and Bible-teaching audio, with downloadable MP3 and slideshow packages: \url{https://globalrecordings.net/en/language/zpx}
\end{itemize}

\section{Curriculum Development Context}
The language is reported to be used as a first language by adults only and is not known to be taught in schools. A curriculum would have to start from new documentation work. The AILLA field notes and any Glottolog-listed references could serve as a starting point. Most of the material found is archival or religious, so it is not suited to classroom use as it stands.

\end{document}